\section{Apple: five statements constrain one year}
\label{sec:apple-statements}

The debt example leaves a larger question: can the reported profit, changes
in equity and cash flows describe the same financial year? Apple's five
consolidated statements supply a useful starting point. Their fiscal 2024
amounts refer to the year ending 28 September 2024; the opening balance sheet
is dated 30 September 2023. All monetary amounts below are millions of US
dollars. The operations statement's per-share amounts and weighted share
counts keep their separate units \citep[printed pp.~29--33]{apple2024}.

Products revenue of 294,866 and services revenue of 96,169 sum to 391,035.
Their costs, 185,233 and 25,119, leave gross margin of 180,683. Research and
development of 31,370 and selling, general and administrative expense of
26,097 reduce that amount to operating income of 123,216. Other income of
269 and an income-tax provision of 29,749 give net income of 93,736
\citep[printed p.~29]{apple2024}. This same net income enters both the equity
statement and the operating-cash reconciliation. It is one period result
appearing in three statements, so it must not be added three times.

Other comprehensive income supplies the next connection. Foreign-currency
translation contributes 395 after tax. Derivative fair-value changes of
$-832$ and reclassification adjustments of $-1{,}337$ contribute $-2{,}169$;
marketable-debt fair-value changes of 5,850 and reclassifications of 204
contribute 6,054. Thus OCI is $395-2{,}169+6{,}054=4{,}280$, and total
comprehensive income is $93{,}736+4{,}280=98{,}016$
\citep[printed p.~30]{apple2024}. Adding the period OCI to opening accumulated
other comprehensive loss of $-11{,}452$ gives the closing $-7{,}172$.
That accumulated balance, rather than the period flow, enters the balance
sheet.

Despite positive comprehensive income, total equity falls. The equity
statement explains the difference: common-stock issuance adds 1,423 and
share-based compensation adds 12,034; withholding reduces the capital
component by 3,993 and retained earnings by 1,612; declarations reduce
retained earnings by 15,218 and repurchases by 95,846
\citep[printed p.~32]{apple2024}. These movements, apart from comprehensive
income, total $-103{,}212$. The complete equity reconciliation is therefore
\[
62{,}146+98{,}016+1{,}423+12{,}034-3{,}993-1{,}612
 -15{,}218-95{,}846=56{,}950.
\]
The components reconcile separately: capital rises from 73,812 to 83,276,
the accumulated deficit moves from $-214$ to $-19{,}154$, and accumulated
other comprehensive loss moves from $-11{,}452$ to $-7{,}172$. These are
the three closing equity components reported on the balance sheet. Their
agreement prevents a model from explaining the correct total with the wrong
allocation among capital, retained earnings and accumulated OCI.

The same change appears in net assets. Table~\ref{tab:apple-statement-join}
compares the reported endpoints; its changes are derived subtractions
\citep[printed p.~31]{apple2024}. Assets rise by 12,397 while liabilities
rise by 17,593, leaving a 5,196 decrease in equity, exactly the decrease
from the equity reconciliation.

\begin{table}[htbp]
\centering
\begin{tabular}{lrrr}
\toprule &30 September 2023&28 September 2024&Change\\\midrule
Assets&352,583&364,980&12,397\\
Liabilities&290,437&308,030&17,593\\
Equity&62,146&56,950&$-5{,}196$\\\bottomrule
\end{tabular}
\caption{The balance-sheet change agrees with the complete equity movement.
The agreement constrains a history without identifying its transactions.}
\label{tab:apple-statement-join}
\end{table}

Cash follows a different reconciliation. Starting from the same net income,
the cash-flow statement adds depreciation and amortisation of 11,445,
share-based compensation expense of 11,688 and other adjustments of
$-2{,}266$. These three adjustments sum to 20,867. Its six operating-asset
and liability adjustments are net receivables $-3{,}788$, vendor non-trade
receivables $-1{,}356$, inventory $-1{,}046$, other assets $-11{,}731$,
payables $+6{,}020$ and other liabilities $+15{,}552$. They sum to 3,651,
giving operating cash of $93{,}736+20{,}867+3{,}651=118{,}254$
\citep[printed p.~33]{apple2024}. The full indirect bridge is thus present;
none of its broad adjustment rows has been replaced by a presumed change in
a similarly named balance-sheet account.

Investing cash is also positive: security purchases of $-48{,}656$,
maturities of 51,211, sales of 11,135, plant-acquisition payments of
$-9{,}447$ and other investing cash of $-1{,}308$ sum to 2,935. Financing
uses 121,983: employee-award tax payments of 5,441, dividend payments of
15,234, share repurchases of 94,949, debt repayments of 9,958 and other
outflows of 361 outweigh net commercial-paper proceeds of 3,960. The source
shows a dash for net term-debt issuance; the displayed total treats that row
as making no contribution, while the source reading retains the dash.
The three activity totals give
\[
118{,}254+2{,}935-121{,}983=-794,
\qquad 30{,}737-794=29{,}943.
\]
The opening and closing amounts here are cash, cash equivalents and restricted
cash and cash equivalents, as defined by the cash-flow statement. Opening
balance-sheet cash is only 29,965, so the cash-flow population exceeds it by
772; both closing amounts are 29,943. Consequently, the balance-sheet line
falls by 22 while the cash-flow total falls by 794. The earlier 772 remains
unallocated in the retained evidence. It cannot be relabelled as the closing
2.6 billion escrow, which Note 4 explicitly includes in reported cash and
equivalents \citep[printed p.~33 and Note 4, printed p.~37]{apple2024}.

These joins expose the missing transaction links more precisely. Annual
receivables rise by 3,902, leaving 114 after adding their $-3{,}788$ indirect
adjustment. Inventory rises by 955, leaving $-91$ after its $-1{,}046$
adjustment. Payables rise by 6,349, exceeding their 6,020 adjustment by 329.
Even the exact opposite vendor-receivable movements identify neither gross
purchases nor collections. The notes examined below must supply the relevant
recognition, valuation, currency and settlement information; the residuals
themselves supply no such explanation.

Nor can the equity rows simply replace financing cash flows. The reported
repurchase reduction of 95,846 exceeds cash repurchases of 94,949 by 897;
withholding reductions total 5,605 against cash tax payments of 5,441; and
equity compensation of 12,034 differs from the 11,688 expense adjustment.
Cash dividends of 15,234 exceed declarations of 15,218 by 16. Those
differences remain unassigned to particular settlements, fees, taxes or
recognition periods \citep[printed pp.~32--33]{apple2024}. The later note
examples retain their own evidence and conditional histories.

A reconstruction must now satisfy the statement relationships together.
Correct total cash cannot excuse wrong equity components, and correct net
assets cannot excuse omitted cash-flow adjustments. The five statements
constrain a common history, but their agreement does not reveal all gross
events, dates, counterparties or accounting treatments. Those unresolved
links determine which historical alternatives can be supported and what
additional facts a conditional future path still needs.
